\chapter{Literature Review}\label{cha:literature_review}

The development of secure IoT architectures spans several domains, from lightweight edge computing to advanced cloud analytics. Recently, the integration of autonomous remediation capabilities has emerged as a critical frontier. This chapter reviews the current state-of-the-art across six areas relevant to the proposed framework: IoT edge anomaly detection, explainable threat classification, AI-driven self-healing networks, privacy-preserving collaborative detection, software-defined and automated network-level remediation, and edge-native model compression. The first three areas correspond directly to the architecture presented in Chapter \ref{cha:methodology}; the latter three are reviewed because they bear directly on the two extensions this report's Future Work section (Chapter 5) proposes-federated collaboration across gateways and on-device small language models-and because, in the case of automated network-level remediation, they surface closely related prior work against which this project's contribution must be positioned. The chapter closes by drawing all six threads together and positioning the proposed framework against the gaps identified.

\section{Anomaly Detection in IoT Edge Environments}

Traditional IDS, such as signature-based detectors, are notoriously ineffective against zero-day vulnerabilities and heavily polymorphic IoT malware. Consequently, researchers have shifted toward anomaly-based detection utilizing Machine Learning (ML). Given the severe computational constraints of IoT devices, recent literature emphasizes decentralized, edge-based detection mechanisms \cite{Gamlo2026}. Deep learning models, particularly Recurrent Neural Networks (RNNs) and Long Short-Term Memory (LSTM) networks, have shown high efficacy in capturing the temporal sequences inherent in network traffic. However, these models are often computationally expensive. GRU autoencoders have been proposed as a more lightweight alternative for capturing temporal dependencies \cite{Yao2023}. Concurrently, traditional one-class classification models, such as the Isolation Forest, are highly regarded for their low overhead and robustness in identifying outliers in high-dimensional data \cite{Liu2008}.

The preference for GRU over LSTM in edge deployments is largely architectural: a GRU merges the input and forget gates of an LSTM cell into a single update gate and removes the separate cell state, reducing the number of trainable parameters and the per-step computation. This makes it comparatively cheaper to run on constrained edge hardware while still retaining the ability to model sequential dependencies in traffic, which is the property this project relies on to build the temporal embedding described in Chapter \ref{cha:methodology}. Isolation Forest, in turn, is attractive for the same resource-constrained setting because it isolates outliers using random recursive partitioning rather than distance or density computations, giving it near-linear scaling with dataset size and no requirement for a labeled attack set during training. This project adopts a related but distinct one-class family for the same reasons of computational lightness and label-free training: Deep Support Vector Data Description (Deep SVDD), which replaces Isolation Forest's random partitioning with a compact neural embedding trained jointly with the anomaly objective, and which Chapter \ref{cha:methodology} motivates specifically by its compatibility with the GRU's own learned temporal embedding-a fusion that a fixed-partition tree method cannot exploit in the same way, since it has no mechanism for co-adapting to a neural feature space during training.

Recent studies have begun exploring hybrid ensembles that combine deep temporal feature extraction with classical anomaly detectors. For instance, combining autoencoders with Isolation Forests has demonstrated improved precision in resource-constrained fog and edge computing environments \cite{Sadaf2020}. However, directly applying these hybrids to IoT edge nodes often results in high false-positive rates due to the non-stationary nature of IoT traffic, since a single-shot anomaly flag cannot distinguish a brief, benign traffic burst from the onset of a sustained attack. The proposed framework advances this paradigm by utilizing a GRU-based temporal embedding fused with a Deep SVDD boundary \cite{Ruff2018}, reinforced by a persistent rolling risk score (Chapter \ref{cha:methodology}) that filters out transient spikes, ensuring high fidelity in sustained threat detection rather than relying on a single anomalous window to trigger a response.

A related and still-open problem in this literature is concept drift: the statistical properties of the traffic used to define ``normal'' behaviour are not fixed but evolve over the lifetime of a deployment, as device firmware updates, changing usage patterns, and new services alter what a benign fingerprint actually looks like. Lu et al., in a widely cited review of the concept-drift literature, distinguish abrupt drift, in which behaviour changes suddenly, from gradual or incremental drift, in which it shifts slowly enough that a detector trained on stale data degrades without an obvious trigger \cite{Lu2018}. A one-class detector trained once on a benign baseline and never revisited is vulnerable to precisely this failure mode: as a deployment ages, the boundary the model learned drifts further from the population it is actually observing, and either false positives accumulate as benign-but-new behaviour is flagged, or the boundary loosens until genuinely malicious traffic falls inside it. The rolling risk score developed in Chapter \ref{cha:methodology} does not solve concept drift outright, since it operates on the output of a fixed ensemble rather than retraining that ensemble online, but its decay mechanism is a partial mitigation in the same spirit: by continuously discounting evidence that is not reinforced by subsequent windows, it prevents a stale or marginally miscalibrated boundary from producing a single persistent false alarm, even though periodic retraining of the underlying GRU--Deep SVDD ensemble remains necessary to address drift at the model level rather than at the scoring level.

\section{Threat Classification and XAI}

While edge nodes are highly effective at detecting binary anomalies, identifying the specific attack vector (e.g., Mirai botnet \cite{Antonakakis2017}, DDoS, cryptojacking) requires the computational power of a centralized or cloud tier. Tree-based ensemble models, particularly Extreme Gradient Boosting (XGBoost) and LightGBM, have consistently outperformed neural networks on tabular network flow datasets (e.g., CICIDS) in terms of accuracy and inference speed \cite{Sama2025}. Botnets such as Mirai are a useful illustration of why fine-grained classification matters operationally: once a device is compromised and enlisted into a botnet, the resulting traffic pattern is behaviorally distinct from, say, a data-exfiltration or reconnaissance attempt, and a coarse binary anomaly flag alone gives the cloud tier no basis for choosing between different healing actions.

Despite their high accuracy, these ML classifiers suffer from a lack of transparency. The ``black-box'' nature of advanced ML models obscures the decision-making process, making it difficult for network administrators to trust automated systems. XAI has emerged to bridge this gap. SHapley Additive exPlanations (SHAP), introduced by Lundberg and Lee, provide local and global interpretability by quantifying the exact, additive contribution of each network feature (e.g., packet length, flow duration) to the final classification \cite{Lundberg2017}. Recent research has heavily advocated for integrating SHAP into IoT security to provide actionable insights to human analysts \cite{Ogunseyi2026}. In the proposed architecture, however, the use of XAI is pivoted: rather than solely aiding human analysts, the SHAP outputs are directly fed into an autonomous multi-agent AI system, providing the necessary forensic context for machine-driven decision-making. This distinction matters because a SHAP vector consumed by a human analyst only needs to be visually legible (e.g., a summary plot), whereas a SHAP vector consumed by an LLM agent must be serialized into a structured, machine-parseable form that preserves feature attribution semantics, which is a requirement this project addresses directly in the cloud-tier design.

The trade-off between detection accuracy and explanation quality, introduced in Chapter \ref{cha:introduction} in the context of analyst trust, is also documented at the technical level in the wider XAI-for-cybersecurity literature. Oriaro and Mishra's systematic review, synthesising 169 peer-reviewed studies on explainable AI in cybersecurity, reports that model-agnostic explainers such as SHAP and Local Interpretable Model-agnostic Explanations (LIME) are the most frequently applied techniques across the surveyed literature, appearing in roughly 40\% and 27\% of the reviewed studies respectively, but that generating these explanations at inference time adds a measured average overhead on the order of 250 milliseconds per alert in large-scale deployments \cite{Oriaro2025}. For a framework such as the one proposed in this report, in which the SHAP explanation stage sits directly in the critical path between cloud-tier classification and the multi-agent core's decision to act, this overhead is a design constraint rather than a footnote: it must remain small relative to the multi-second remediation-latency budget reported in Chapter 4 for the explanation step not to become the dominant contributor to end-to-end response time. The same review also documents a security-specific risk that is easy to overlook when XAI is treated purely as a usability feature: adversarial manipulation of the explanation itself, in which an attacker crafts inputs designed to produce a benign-looking SHAP or LIME attribution regardless of the underlying model's true classification, was identified as a demonstrated attack vector in approximately 15\% of the adversarial-robustness studies the review covers \cite{Oriaro2025}. This risk is a further argument, alongside the trust-erosion argument developed in Chapter \ref{cha:introduction}, for why this project treats the SHAP output as one input among several-alongside the XGBoost classification label and the raw feature vector-that the Decision and Escalation Agents reason over when consulted, rather than as a sole and unquestioned ground truth.

\section{Agentic AI, Multi-Agent Systems, and Self-Healing Networks}

The traditional approach to IoT fault management and incident response is overwhelmingly manual and reactive. As IoT networks scale, human-in-the-loop remediation introduces unacceptable latency. The concept of ``self-healing'' networks-systems capable of autonomous fault detection, isolation, and service restoration-has gained significant traction \cite{Manju2025}.

Historically, self-healing frameworks relied on rigid, rule-based expert systems or static Software-Defined Networking (SDN) policies. Such rule-based systems are precise but brittle: they can only react to conditions their designers explicitly anticipated, and extending them to a new attack type or network topology requires manual reprogramming. The advent of LLMs has catalyzed the development of ``Agentic AI'' in network security \cite{Nageshwaran2026}, as an alternative that can reason over unstructured evidence (e.g., a forensic report) rather than only matching against predefined rules. Unlike traditional chatbots, autonomous LLM agents can maintain state, invoke external tools, and execute complex workflows. Multi-Agent Systems (MAS) distribute cognitive workloads among specialized agents-for instance, separating data analysis, compliance checking, and action orchestration-thereby improving system reliability and reducing hallucination risks compared to a single monolithic agent attempting every step of the pipeline \cite{Vinay2025}.

A critical challenge in applying LLM agents to network administration is ensuring that the generated commands respect physical hardware limitations and network topologies; an ungrounded LLM may propose a remediation action that is technically fluent but operationally invalid (e.g., isolating a VLAN that does not exist on the target device). Recent approaches utilize RAG to ground LLM outputs in verified documentation and system constraints, retrieving only the passages relevant to the device and incident at hand before the LLM commits to an action \cite{Lewis2020}. While early agentic frameworks have been proposed for general IT operations, few studies have successfully closed the loop in constrained IoT environments, where actions must additionally be verified against a live, evolving risk signal rather than assumed successful once dispatched. The proposed framework addresses this gap by orchestrating a MAS that digests SHAP-enriched threat data, queries a RAG system for verified mitigation protocols, and autonomously enforces healing actions (such as dynamic VLAN isolation and bandwidth reallocation) directly at the edge layer, with the rolling risk score itself serving as the closed-loop success signal.

The specific combination of LLM-based reasoning with network-level enforcement pursued here is not unprecedented, and recent work illustrates both the promise and the limitations of closely related approaches. Swileh and Zhang propose a zero-training detection and mitigation scheme for decentralised software-defined networks in which a large language model classifies port-level traffic statistics through in-context learning alone, then triggers mitigation directly at the offending switch port, reporting near-perfect detection accuracy under DDoS conditions without any model fine-tuning \cite{Swileh2025}. A follow-up study extends this approach to Carpet-Bombing DDoS attacks-a low-and-slow variant that spreads traffic across many victim subnets specifically to evade volumetric thresholds-by adding a retrieval-augmented layer that grounds the LLM's classification in semantically retrieved traffic patterns rather than in-context examples alone \cite{Swileh2026}. Separately, Islam and Saha propose a two-stage pipeline that pairs a deep-learning classifier with a RAG system that retrieves from authoritative cybersecurity knowledge sources, including NIST and MITRE ATT\&CK material, to generate structured, citation-grounded mitigation reports for network intrusions \cite{IslamSaha2026}.

These three systems demonstrate that LLM-driven classification and RAG-grounded remediation guidance are each independently viable, and each was evaluated in isolation from the constraints this project must satisfy simultaneously. None of the three couples its language-model reasoning to a persistent, decaying risk signal of the kind developed in Chapter \ref{cha:methodology}; each instead issues or recommends an enforcement action on the strength of a single classification decision, which reintroduces at the enforcement stage the same single-shot limitation already identified in Section 2.1 for anomaly detection. Nor does any of the three target the severely resource-constrained edge tier that motivates pushing initial detection off the cloud path entirely, as argued in Section 1.1: all three operate on SDN switches or general-purpose network infrastructure with comparatively abundant compute, rather than on Class 0 to Class 2 IoT endpoints and their gateways. The proposed framework can therefore be read as extending this emerging line of LLM-grounded network defence with the two properties it currently lacks: a stateful detection signal that persists and decays across windows rather than a per-flow classification, and an edge-cloud division of labour suited to constrained IoT endpoints rather than to comparatively unconstrained network infrastructure.

\section[Federated Intrusion Detection]{Federated Learning for Privacy-Preserving\\Collaborative Intrusion Detection}

The cascaded hybrid ensemble described in Chapter \ref{cha:methodology} is, in its baseline form, trained and evaluated at a single gateway; each edge node learns its own model of benign behaviour from its own traffic. This is appropriate for the detection problem this project addresses directly, but it leaves open a question that this report's Future Work section flags explicitly: how should a fleet of many such gateways, deployed across different sites and operators, share what each has learned about emerging attack patterns without exposing the raw traffic-and, by extension, potentially sensitive operational or patient data-that each has observed. Federated Learning (FL), introduced by McMahan et al. as a method for training a shared model across many clients by exchanging model updates rather than raw data \cite{McMahan2017}, is the dominant paradigm the literature has converged on for this problem in the intrusion-detection context specifically.

A recent systematic review by Hernandez-Ramos et al., synthesising more than 100 FL-enabled IDS studies, proposes a taxonomy spanning the choice of local ML model, the aggregation function used to combine client updates, the datasets used for evaluation, and the FL framework or library used for implementation \cite{HernandezRamos2025}. The review's central finding is that FL-based IDS research has grown rapidly, but that most proposed systems remain evaluated under idealised conditions-balanced data, homogeneous client hardware, and a benign aggregation server-that are not representative of a real IoT deployment \cite{HernandezRamos2025}. A complementary survey by Zhang et al. structures the same literature around six stages of an FL-IDS pipeline (data preparation, local training, aggregation strategy, communication, evaluation and deployment) and identifies high communication latency, elevated false-positive rates under non-independent-and-identically-distributed (non-IID) client data, and vulnerability to poisoning attacks from malicious participants as the three most persistent open problems across the surveyed work \cite{Zhang2025}.

Each of these three problems maps directly onto a characteristic of the endpoint population described in Section 1.1. Communication latency is aggravated, not merely present, in a fleet of Class 0 to Class 2 devices whose gateways already operate under the bandwidth constraints discussed in Section 1.1.2; a federated round that requires every gateway to transmit a full model update on a fixed schedule competes with the same operational traffic that motivated pushing detection to the edge in the first place. Non-IID data is close to guaranteed rather than an edge case: a hospital's infusion-pump traffic and a water utility's SCADA telemetry, both belonging to the endpoint population this project addresses, have almost nothing in common at the traffic-feature level, so any federation across sites must contend with genuinely heterogeneous local distributions rather than statistical noise around a shared mean. Poisoning risk is compounded by the same repeated, targeted-attacker profile documented in Section 1.1.3: an adversary who has already compromised one gateway in a federation has an incentive to submit corrupted updates specifically to degrade the shared model's ability to detect the attack pattern that compromised it, which is a materially different threat model from the accidental noise that early FL aggregation methods were designed to tolerate.

Beyond model accuracy, FL for IDS also raises privacy and integrity concerns that a systematic review by Mothukuri et al. treats as a subject in its own right, distinguishing threats to the confidentiality of client updates, which can under some conditions be inverted to reconstruct approximations of local training data, from threats to model integrity, in which malicious clients submit crafted updates to bias the global model \cite{Mothukuri2021}. Secure aggregation protocols and differential-privacy noise injection are the two dominant mitigations the review identifies for the former; robust aggregation rules that discount statistically anomalous client updates are the dominant mitigation for the latter \cite{Mothukuri2021}. Both mitigations impose an additional computational or communication cost on top of the baseline FL round, which compounds rather than replaces the bandwidth constraint already discussed in Section 1.1.2, and is a further reason the FL extension proposed as future work is treated as requiring dedicated design work rather than a direct application of an off-the-shelf FL library.

These findings do not change the architecture presented in Chapter \ref{cha:methodology}, which is deliberately scoped to single-gateway detection and risk scoring; they instead substantiate, with reference to the current state of the literature, why the Federated Learning extension proposed in this report's Future Work section is treated as a distinct and non-trivial undertaking rather than a routine engineering extension. Collaboratively updating the GRU--Deep SVDD ensemble across gateways would require resolving the same non-IID and poisoning-robustness questions that the reviewed literature identifies as unresolved, in addition to the communication-budget constraints specific to the endpoint population this project targets.

\section{Software-Defined Networking and Automated Network-Level Remediation}

The healing actions this project's orchestration agent is designed to issue-VLAN isolation, bandwidth reallocation, and port blocking-are not novel primitives; they are the standard enforcement vocabulary of Software-Defined Networking (SDN), in which a logically centralised controller can reprogram forwarding behaviour across a network without manual reconfiguration of individual switches. Jain et al., in a comprehensive survey of DDoS detection and mitigation strategies for SDN environments, catalogue the enforcement actions available once an attack is detected-traffic rate limiting, flow-rule installation to drop or redirect malicious flows, and dynamic access-control list updates among them-and note that the SDN control plane's programmability is precisely what makes automated, sub-second remediation possible in principle, in contrast to the manual router and switch reconfiguration that conventional networks require \cite{Jain2024}. The same survey observes, however, that the majority of SDN-based mitigation systems in the literature couple this programmability to a single-shot detection decision: a flow is classified as malicious or benign once, and the corresponding enforcement action is applied without a subsequent verification step to confirm that the action actually suppressed the attack \cite{Jain2024}. This recurs as the same single-shot limitation already identified in Section 2.1 for anomaly detection, but at the enforcement stage rather than the detection stage: an SDN controller that blocks a flow because a classifier once flagged it has no built-in mechanism to notice that the block failed, was circumvented, or was unnecessary, and few of the surveyed systems close this loop.

The same survey catalogues the datasets and simulation tools most commonly used to evaluate SDN mitigation systems-among them Mininet-based topologies paired with synthetic attack-traffic generators and flow datasets adapted from conventional network security research-and observes that comparatively few of the reviewed systems are evaluated against attack traffic drawn from real IoT botnets rather than generic DDoS generators \cite{Jain2024}. This gap is significant for the present project, since Section 1.1.3 established that the endpoint population motivating this report generates traffic with device-specific characteristics-the constrained protocol stacks and periodic reporting patterns typical of Class 0 to Class 2 devices-that differ systematically from the general-purpose network traffic most SDN mitigation research is evaluated against; an enforcement mechanism validated only on conventional DDoS traffic offers no guarantee that it will behave correctly against the narrower, more repetitive traffic signatures of an IoT botnet.

A second relevant strand of this literature concerns the specific enforcement primitive of dynamic VLAN and access segmentation as a containment mechanism, rather than as a purely static network-design choice. Isolating a compromised or suspect endpoint to its own VLAN, or throttling the bandwidth available to a specific port, has the property that-unlike a flow-level firewall rule, which can be circumvented by an attacker who alters flow characteristics-it acts on a hardware-level forwarding boundary that a compromised device typically cannot bypass through application-layer evasion alone. This is precisely why the framework proposed in this report favours these primitives over content-based filtering: the objective is not to inspect and drop specific malicious payloads, which a sufficiently adaptive attacker can reformulate, but to physically constrain what a compromised device is capable of reaching on the network until its risk score decays, which is a property of the segmentation boundary rather than of the traffic content.

The most direct precedent for combining these SDN enforcement primitives with language-model-driven decision-making is the body of work already discussed in Section 2.3 in the context of agentic AI-Swileh and Zhang's port-level LLM classifier \cite{Swileh2025}, its Carpet-Bombing-specific RAG extension \cite{Swileh2026}, and Islam and Saha's RAG-grounded mitigation-report generator \cite{IslamSaha2026}-each of which issues or recommends an SDN-style enforcement action once its classification stage fires. What distinguishes the proposed framework from this precedent, as argued above, is not the enforcement vocabulary itself, which is shared, but the fact that enforcement here is verified against a continuing risk signal rather than treated as a terminal action, and that the detection stage feeding that enforcement decision is explicitly designed for a gateway-class device rather than for an SDN switch with a comparatively unconstrained control plane.

\section{Edge-Native Model Compression and On-Device Inference}

A second direction flagged in this report's Future Work section is the eventual replacement of the cloud-hosted multi-agent LLM core with a heavily quantised, domain-specific small language model (SLM) running directly on edge-gateway hardware, reducing the framework's dependence on a continuously available cloud connection. This ambition sits within a broader and increasingly mature body of work on TinyML: the practice of compressing deep learning models, through quantisation, pruning, knowledge distillation, and architecture search, until they fit within the memory and power budget of embedded or near-edge hardware. Rajapakse et al., surveying what they term ``reformable'' TinyML systems, distinguish static compression, in which a model is compressed once before deployment and never changes, from reformable approaches, in which a deployed model can be further adapted, retrained, or reconfigured in the field as conditions change \cite{Rajapakse2023}. This distinction is directly relevant to the design tension already identified in Section 1.1.2 between fixed on-device capability and the need to respond to an evolving threat landscape: a statically compressed detector, once deployed to a gateway, faces exactly the concept-drift problem discussed in Section 2.1, whereas a reformable one retains some capacity to adapt without a full cloud-mediated retraining cycle.

Heydari and Mahmoud, in a survey focused specifically on on-device inference, report that the dominant obstacle to deploying transformer-based language models-as opposed to the convolutional or recurrent architectures more commonly targeted by earlier TinyML work-on constrained hardware is not raw parameter count alone but the memory bandwidth required to move activations and attention weights during inference, which does not shrink proportionally with aggressive quantisation the way parameter storage does \cite{Heydari2025}. This has a direct bearing on the feasibility of the SLM-on-gateway direction proposed as future work: even a heavily quantised language model small enough to fit in a gateway's flash storage may still be constrained by the memory-bandwidth ceiling that Heydari and Mahmoud identify as the harder of the two bottlenecks, meaning that model size alone is not a sufficient design target. The same survey notes that current on-device transformer deployments in adjacent domains typically operate at a substantially reduced context length and vocabulary size relative to their cloud-hosted counterparts, trading generality for footprint \cite{Heydari2025}-a trade-off that would need to be revisited carefully for this project's use case, since the multi-agent core's Analytical and Reporting Engines currently depend on a comparatively rich vocabulary of network-security terminology and forensic narrative structure that a heavily reduced vocabulary could constrain.

Positioned against this literature, the current architecture's choice to keep the GRU autoencoder and Deep SVDD boundary at the edge while reserving the LLM-based multi-agent core for the cloud tier is consistent with where the TinyML literature places the achievable frontier today: lightweight numerical and sequence models are comparatively mature targets for edge compression, as reflected in the GRU-over-LSTM design choice discussed in Section 2.1, while compressed language models capable of the semantic reasoning this project's Decision and Escalation Agents perform remain, per Heydari and Mahmoud's survey, at an earlier and less settled stage of feasibility \cite{Heydari2025}. This is offered as the literature-grounded justification for treating edge-native LLM deployment as future work rather than as part of the present system's scope.

\section{Research Gap and Positioning of the Proposed Framework}

The six areas reviewed above serve two distinct purposes. Anomaly detection, threat classification and XAI, and agentic AI and self-healing networks (Sections 2.1--2.3) describe the state of the art that the architecture in Chapter \ref{cha:methodology} directly builds on and departs from. Federated learning, SDN-based remediation, and edge-native model compression (Sections 2.4--2.6) describe adjacent literatures that the present system does not yet incorporate but that the Future Work section of this report explicitly identifies as extensions, and reviewing them here grounds those extensions in the current state of their respective fields rather than treating them as unexamined possibilities. Drawing the first three review threads together against the limitations identified in Section \ref{cha:introduction} (Problem Statement), three concrete gaps motivate this project:

\begin{enumerate}
	\item \textbf{Single-shot vs. persistent detection:} existing hybrid anomaly detectors improve precision over standalone models but still commit to a decision on a single traffic window, leaving them exposed to transient bursts of legitimate IoT traffic. The proposed two-stage rolling risk score (Chapter \ref{cha:methodology}) treats detection as an accumulating, decaying signal rather than a one-shot classification, so that only sustained deviations escalate to the cloud tier.
	\item \textbf{Human-only vs. machine-consumable XAI:} SHAP and similar techniques have been validated as tools for human analysts, but the literature stops short of using their outputs as a direct input to an automated decision-making system. This project extends SHAP's role from visualization to a machine-to-machine interface that an LLM agent can reason over directly.
	\item \textbf{Diagnostic vs. restorative automation:} agentic AI and RAG grounding have been explored for general IT operations, but rarely combined into a closed loop that both issues a remediation command and verifies its effect using a live, device-specific signal in a resource-constrained IoT setting. The proposed orchestrator agent closes this loop by monitoring the same rolling risk score used for detection, iterating its strategy or escalating to a human operator when the score fails to decay.
\end{enumerate}

These three gaps directly correspond to several of the five objectives stated in Section \ref{cha:introduction} (Objectives and Scope)-persistent detection and machine-consumable XAI map onto the objectives concerning fast, resilient response and explainable, trustworthy decision-making respectively, while the closed-loop diagnostic-vs-restorative gap underpins both the layered-intelligence objective and the self-directed policy-learning objective-and the remainder of this report is organized around addressing them.